离线评估 AUC 0.88，上线之后 0.74。

排查一圈：没有分布漂移，没有代码缺陷，训练数据也没问题。最后发现是标准化在划分之前做的——验证折的均值与方差泄漏进了训练。把整条流程搬进每一折内部重做，离线指标掉到 0.76，与线上对上了。

这个故事里没有一步涉及新算法。出问题的是\emph{关于未见数据的断言怎样才算数}，而这恰好是机器学习真正的主题。模型族、损失函数、优化器都是可替换的零件；不可替换的是那句断言的成立条件。

前面十二个部分已经把这些条件分头讲过：线性代数说清了投影与低秩，概率论说清了条件独立与依赖，凸优化说清了什么时候局部即全局，数理统计说清了置信度属于程序，信息论说清了边界在哪里。这一部分做的事，是把它们\textbf{按一个具体问题重新排列}——从一堆数据到一个可以依赖的判断，中间要经过哪些环节，每个环节由谁管辖。

\begin{center}
\begin{tikzpicture}[x=1cm,y=1cm,font=\small,>=stealth]
  \draw[black!45,fill=blue!8] (0,3.3) rectangle (3.9,4.3);
  \node[align=center] at (1.95,3.8) {问题翻译};
  \draw[black!45,fill=blue!8] (0,1.95) rectangle (3.9,2.95);
  \node[align=center] at (1.95,2.45) {模型族与归纳偏置};
  \draw[black!45,fill=blue!8] (0,0.6) rectangle (3.9,1.6);
  \node[align=center] at (1.95,1.1) {推断与训练};
  \draw[black!45,fill=blue!8] (0,-0.75) rectangle (3.9,0.25);
  \node[align=center] at (1.95,-0.25) {评估与决策};
  \foreach \y in {3.3,1.95,0.6} \draw[->,black!55] (1.95,\y) -- (1.95,\y-0.35);
  \node[align=left,anchor=west,font=\footnotesize,black!70] at (4.5,3.8)
    {目标写歪，后面全白做（第一章）};
  \node[align=left,anchor=west,font=\footnotesize,black!70] at (4.5,2.45)
    {线性 / 树 / 核 / 概率图，各自假设了什么（第二至九章）};
  \node[align=left,anchor=west,font=\footnotesize,black!70] at (4.5,1.1)
    {写得出却算不出时，采样还是变分（第八章）};
  \node[align=left,anchor=west,font=\footnotesize,black!70] at (4.5,-0.25)
    {泛化、因果、以及信息本身的上限（第十至十三章）};
\end{tikzpicture}

\emph{十三章按这条纵线排列。往下走一格，能做的承诺就多一层，需要的假设也多一组——两者永远成对出现。}
\end{center}

\subsection{五条贯穿始终的判断}

\textbf{最难的一步是翻译，不是拟合。}把业务诉求写成决策变量、目标函数与约束，这一步错了，后面解得再准也是错的。只优化点击总数会催生夸张标题，只考核平均处理时长会催生提前结单——目标函数会被精确地钻空子，这与凸优化里那条判断是同一件事。

\textbf{泛化是关于未见数据的断言，所以评估必须真的模拟``未见''。}任何在划分之前动过全量数据的步骤都会破坏它，开头那个故事就是最常见的形态。一个便宜的自检是把标签打乱重跑，指标应当回到随机水平；回不去就说明有泄漏。

\textbf{每个模型族都是一组关于世界形状的假设。}线性模型假设效应可加，树假设决策边界轴对齐，核方法把假设写进核函数，概率图把假设写成``哪些依赖不存在''。选模型就是在选假设，而假设写错了不会报错，只会让结论悄悄有偏。所以该问的不是``哪个模型更强''，是``哪组假设更接近我这份数据''。

\textbf{三类误差必须分开记账。}训练损失降不下去是优化问题，训练低而测试高是统计问题，测试好而线上差通常是分布或口径问题。混淆它们会导致该加数据的时候去调优化器。相应地，统计学习理论给出的界极其宽松，它的用处是指出哪些因素影响泛化，而不是预测泛化误差的数值。

\textbf{预测世界不等于改变世界。}相关性由联合分布决定，因果性还需要额外的结构假设——同一个联合分布可以对应多张因果图，观察数据无法在它们之间区分。控制变量也不是越多越好：控制混杂消除偏差，控制碰撞点制造偏差。任何``改了这个会怎样''的问题，都已经越出了预测的范围。

\subsection{十三章怎么安排}

\hyperref[ch:109]{``学习问题与可信评价：先确定究竟要学什么''}是地基，几乎不含算法却决定后面所有话该怎么说。\hyperref[ch:110]{``线性模型：先让假设透明，再让模型复杂''}紧随其后——它至今不被取代的理由不是表达力，是系数有含义、不确定性可量化、失效方式已知。

想按模型族取用，\hyperref[ch:111]{``树、集成与邻域：用局部结构代替全局公式''}、\hyperref[ch:117]{``核方法与 Gaussian process：把函数当作随机对象''}、\hyperref[ch:114]{``概率模型与隐变量：从``预测标签''到``解释数据怎样产生''''}、\hyperref[ch:115]{``概率图与序列模型：用条件独立压缩联合分布''}各自独立成篇，共同点是每一章都先交代自己假设了什么。\hyperref[ch:112]{``降维与低秩建模：数据为何看似高维''}与\hyperref[ch:113]{``聚类与数据几何：没有标签时结构从何而来''}处理没有标签时的结构问题，两章都要记住：高维下的谱会系统性失真，肘部可能是噪声造的。

\hyperref[ch:116]{``采样与近似推断：当目标量写得出却算不出''}是贝叶斯方法与生成模型的共同入口。\hyperref[ch:118]{``统计学习理论：训练误差对真风险有多乐观''}与\hyperref[ch:119]{``信息论学习：把预测、约束、复杂度与表示放到同一尺度''}给出两类上限——前者关于假设类容量，后者关于信息本身；后者的 Fano 界尤其实用，它区分``模型还不够好''与``特征里根本没有这个信息''。

\hyperref[ch:120]{``因果推断：预测世界不等于改变世界''}建议完整读完，即使暂时不做因果分析——它讲的控制变量陷阱在任何回归解释里都会遇到。\hyperref[ch:121]{``统计物理与自由能：平均场近似的来源与边界''}是选读，价值在于把变分推断的自由能视角讲清，同时诚实划出哪些物理类比只是类比。

这一部分与下一部分的分工是这样的：这里的模型族大多把``用什么表示''交给人来设计，而\hyperref[ch:122]{``深度学习：让表示、梯度与数据结构在复合函数中相遇''}让表示本身成为可优化的对象。代价是几乎所有保证都退化为经验规律——那一部分开头会把这笔账摊开讲。
